Rebuttal (Reviewer 1)

We thank the reviewer for the careful reading and highly constructive feedback. The comments helped us identify several places where notation and exposition can be improved. Below we provide point-by-point responses and the corresponding camera-ready changes.

Q1 (Figure clarity / momentum comparisons):  
We agree and will redraw Figure 2 with a cleaner two-color convention and matched full-rank vs low-rank curves at fixed momentum/hyperparameters. Cross-momentum sweeps will be moved to appendix/supplementary. We will also state that these curves diagnose optimization/loss-landscape effects under rank constraints; in our setup, plain SGD (no momentum) can be competitive with momentum variants, consistent with the zero-momentum mean-field gradient-flow view.

Q2 (Citation style consistency):  
We will standardize citation style throughout the manuscript (consistently either Author-Year or numeric style).

Q3 (Meaning of supp(W^0(C1)) / frozen random features):  
We will replace this unclear wording with a standard random-feature richness assumption: frozen first-layer features induce a dense span in an appropriate function class (e.g., L2(mu_X)). We will also separate assumptions by depth: richness at layer 0, and bounded mixing for deeper layers. Concretely, for ell > 1, we assume bounded mixing with ||W^(ell)||_{infty,1} <= rK, which controls channel aggregation and stability.

Q4 (Mean-field width limit):  
This was a wording typo. We mean the infinite-width limit under mean-field parameterization (not NTK/lazy scaling and not muP). We will rephrase this explicitly and state which widths diverge.

Q5 (Matrix-form suggestion in Sec. 3.1):  
We agree that the compact matrix view is useful and will add it in the main text as a complement. We keep the channelized representation as primary because it explicitly exposes the r channel functions f_k, which are the state variables in the mean-field dynamics and the quantities tracked experimentally.

Q6 (Neuron-index notation / expectations over indices):  
We use neuronal embedding in the standard probabilistic sense: each layer is defined on an index probability space (Omega_i, rho_i), expectations are integrals under rho_i, and finite-width neurons are i.i.d. draws from rho_i. Under mean-field scaling, empirical averages converge to these expectations. We will remove ambiguous wording (“indices become continuous”) and add the equivalent pushforward-measure interpretation on Euclidean weight space.

Q7 (Loss definition and gradient-flow structure):  
We will explicitly define the population loss and residual term in Section 3, and state that the mean-field ODE is the gradient flow with schedules xi_ell(t). These definitions are currently in the appendix; in the camera-ready, we will also include concise versions in the main body so the section is self-contained.

Q8 (Learning-rate schedules):  
We will define xi_ell(t) directly in the ODE and state the Euler correspondence W_{n+1} = W_n - epsilon xi_ell(epsilon n) drift(epsilon n), with mean-field limit epsilon -> 0 and standard width scaling.

Q9 (“Better than null” assumption):  
We will restate this as a non-degeneracy condition: initialization must perform strictly better than the trivial constant predictor.

Q10 (ReLU, high probability in r, and assumptions):  
We will clarify that r is the bottleneck rank (number of channels) and that plain ReLU is excluded from the strict non-vanishing-derivative assumption used in the formal implication “vanishing conditional gradient => vanishing residual.” For ReLU, the high-probability-in-r statement is presented as intuition only: under symmetric mixing and non-degenerate channels, full dead-gate alignment is heuristically exponentially unlikely in r (equivalently on an exponentially small input subset where many weights receive zero gradient through dead gates). We also mention sign-structured/NMF-style initialization as design intuition to bias H2 positive at initialization and reduce dead-gate behavior. This sits alongside the following point on randomness: the proofs ask for bounded frozen mixing and feature richness, not a specific measure---finite-support draws for frozen W suffice; unbounded Gaussians in code are a practical convenience. We have not run a dedicated multi-draw variance study over independent frozen-feature realizations at fixed optimization; see the merged reply “Sensitivity to the frozen feature draw” in the main rebuttal document.

Q11 (Meaning of mixing):  
We will define mixing explicitly: W^(ell) recombines r channels into neuron pre-activations, with structural bound ||W^(ell)||_{infty,1} <= rK. We will also state that W^(ell) may be drawn from any bounded-support distribution (or any law satisfying the same almost-sure bound).

Q12 (Notation in Sec. 3.5–3.6):  
We agree the notation was confusing. In this part, the trainable object is factor A, while the prefactor is a norm of the frozen mixing matrix. We will correct the inequality, explicitly separate trainable vs. frozen quantities, and harmonize notation in the camera-ready to avoid switching between W/L symbols for the same mixing entries.

Q13 (Meaning of solution operator F):  
We will define F as the Picard integral map whose fixed point is the ODE solution trajectory, and state this clearly where first introduced.

Q14 (Minor punctuation):  
We will correct the sentence punctuation and improve readability.

Q15 (Theorem 3.2 scope in depth/width):  
The result is intended for arbitrary depth in RF-LR. We will rewrite the theorem statement to specify which widths tend to infinity (all hidden-layer widths under mean-field scaling) and separate this general statement from the depth-2 notation used for the quantitative bound.

Q16 (“High level proof idea” opacity):  
We will either remove the current high-level paragraph from the main text or rewrite it to mirror the formal proof structure in the appendix (dense-span step, stationarity, conditional expectation, and loss-optimality conclusion).

Q17 (“upstream x local” wording):  
We will replace this by standard terminology (“backpropagated signal x local derivative”) and provide explicit definitions.

Q18 (Meaning of partial_2 L):  
We will explicitly define partial_2 L(y, yhat) as the derivative with respect to the second argument (prediction).

Q19 (Meaning of “shifts” and heterogeneous shifts):  
We will define “shift” as neuron-independent additive translation in collapse regimes of full-rank i.i.d. settings, then contrast with RF-LR where frozen mixing induces neuron-dependent channel combinations, yielding heterogeneous effective updates and preventing the “all neurons move identically” mode.

Q20 (Init index notation in Theorem 3.3):  
We will rewrite the notation in standard finite-width terms: n1,n2 are finite layer widths, and Init is a width-indexed initialization family. This will remove ambiguity.

Q21 (“coupling procedure” and distance D):  
We agree the phrase was misleading. We mean the standard finite-width/mean-field embedding comparison (empirical averages vs. expectations under i.i.d. sampled neuron labels), not an additional algorithm. We will define the discrepancy metric explicitly (trajectory-level discrepancy with O(1/sqrt(n_min) + sqrt(epsilon)) scaling, up to the Gronwall factor) and clarify the transport/Wasserstein interpretation.
